\documentclass[aspectratio=169,10pt]{beamer}
\usepackage[utf8]{inputenc}
\usepackage[T1]{fontenc}
\usepackage[english]{babel}
\babelprovide[import,main]{french}
\usepackage{lmodern,amsmath,amssymb,booktabs,graphicx}
\newcommand{\DataCommit}{867144db98eba3c0f5fe7cece4261918a46f1c31}
\newcommand{\DataURL}[1]{https://github.com/halhoulmhamed-droid/inverse-laplace-learning/blob/867144db98eba3c0f5fe7cece4261918a46f1c31/studies/study3/#1}

\definecolor{ink}{RGB}{25,48,73}
\definecolor{accent}{RGB}{23,116,126}
\definecolor{caution}{RGB}{159,76,35}
\setbeamercolor{normal text}{fg=ink,bg=white}
\setbeamercolor{frametitle}{fg=ink}
\setbeamercolor{structure}{fg=accent}
\setbeamercolor{block title}{bg=accent!12,fg=ink}
\setbeamercolor{block body}{bg=accent!4,fg=ink}
\setbeamercolor{alerted text}{fg=caution}
\setbeamertemplate{navigation symbols}{}
\setbeamertemplate{footline}{\hfill\scriptsize Halhoul Mhamed\quad Étude 3\quad\insertframenumber/\inserttotalframenumber\hspace{5mm}\vspace{3mm}}
\setbeamersize{text margin left=9mm,text margin right=9mm}
\hypersetup{pdftitle={Étude 3 : Régularisation des covariances apprises pour une reconstruction robuste de données de Laplace},pdfauthor={Halhoul Mhamed}}
\newcommand{\src}[2]{\href{\DataURL{#1}}{#2}}
\newcommand{\figsource}[1]{\par\tiny\src{figures/#1.npz}{Données}\,|\,\src{src/figures.py}{Script}\,|\,\src{results/final/configuration.json}{Configuration}\,|\,\src{figures/#1.pdf}{Version scientifique}\par}
\newcommand{\plot}[2]{\includegraphics[width=#2\linewidth,height=.61\textheight,keepaspectratio]{../figures/#1.pdf}}
\begin{document}

\begin{frame}[plain]
\begin{columns}[T]
\column{.48\textwidth}\includegraphics[width=\linewidth,height=12mm,keepaspectratio]{../../../assets/logos/university.png}
\column{.48\textwidth}\raggedleft\includegraphics[width=\linewidth,height=12mm,keepaspectratio]{../../../assets/logos/faculty.png}
\end{columns}
\centering
{\scriptsize Université Abdelmalek Essaâdi --- Faculté Polydisciplinaire de Larache\par}
\par\vspace{2mm}
{\large Projet scientifique de Master --- Étude 3\par}
\par\vspace{3mm}
{\LARGE\bfseries Régularisation des covariances apprises\par}
{\large pour une reconstruction robuste de données de Laplace\par}
\par\vspace{4mm}
{\large Halhoul Mhamed\par}
{\scriptsize Master « Modélisation mathématique et applications et apprentissages » --- M2 en cours\par}
\par\vspace{4mm}
\textbf{Professeur encadrant : Pr. Chakir Tajani}\par
\vfill
{\tiny 11 octobre 2026}
\end{frame}

\begin{frame}{Une amélioration ponctuelle, un gain utile non établi}
\begin{block}{Question}
Une covariance apprise régularisée vers une cible H1 peut-elle réduire l'erreur avec tous les capteurs, tout en préservant la reconstruction avec des capteurs partiels ?
\end{block}
\par\vspace{2mm}
\begin{itemize}
\item Le ratio de risque S/M0 observé vaut \textbf{0,9607 à 0,9823} dans le scénario principal.
\item Aucune des cinq bornes supérieures n'atteint le seuil de gain utile fixé à \textbf{0,97} : \alert{conclusion principale non établie}.
\item La marge de non-dégradation de \textbf{3\,\%} sur 32 capteurs contigus est respectée par les cinq bornes conditionnelles.
\end{itemize}
\par\vspace{2mm}
\small Des dégradations apparaissent sous rugosité ou erreur de hauteur. Aucun résultat positif ni aucune supériorité générale n'est présumé.
\end{frame}

\begin{frame}{Un nouveau test, sur un modèle physique inchangé}
\[
q(x)=\sum_{k=0}^{127}c_k\phi_k(x),\qquad
y=M(B_Hc+\sigma\varepsilon)/8,\quad \varepsilon\sim\mathcal N(0,I_{64}).
\]
\begin{itemize}
\item Rectangle harmonique, $H=0,25$ ; 64 moyennes de capteurs. Générateur analytique à 128 modes ; inverse FD à 64 modes.
\item Tous les nouveaux estimateurs prédisent les 64 coefficients : aucune troncature singulière à rang 12.
\item Nouvelle racine de graines \texttt{2026101103}. Apprentissage, validation, pilote et test séparés ; mêmes bruits entre méthodes.
\end{itemize}
\begin{block}{Interprétation des covariances}
Les coefficients cosinus du générateur restent indépendants. Des corrélations en coordonnées singulières peuvent venir du changement de base ; elles ne démontrent pas des corrélations physiques des coefficients.
\end{block}
\end{frame}

\begin{frame}{Mélanger des covariances de signal à énergie égale}
\small
\[
C_{\rm emp}=\frac1{N-1}\sum_i(c_i-\mu)(c_i-\mu)^\top,\quad
Q_{kk}=1+(\pi k)^2,\quad
C_{H1}=sQ^{-1},\quad s=\frac{\operatorname{tr}C_{\rm emp}}{\operatorname{tr}Q^{-1}}.
\]
\[
\boxed{C_\gamma=(1-\gamma)C_{\rm emp}+\gamma C_{H1}},\qquad
\operatorname{tr}C_\gamma=\operatorname{tr}C_{\rm emp}.
\]
\[
\widehat c=\mu+C_\gamma A^\top(AC_\gamma A^\top+\beta I)^{-1}(y-A\mu),
\qquad \beta=\alpha\sigma^2/64>0.
\]
\begin{itemize}
\item $C_\gamma$ est PSD ; le système d'observation est SPD, même si $\gamma=0$ et $N<64$.
\item $\gamma=1$ : H1 affine exact, pénalité $\eta Q$ avec $\eta=\beta/s$, centre $A\mu$.
\item La cible H1 diffère de la vraie covariance nominale. Positivité et solvabilité ne garantissent pas une baisse du risque.
\end{itemize}
\end{frame}

\begin{frame}{Six familles : centres et budgets déclarés}
\small
\begin{center}
\begin{tabular}{llll}
\toprule
Méthode & Construction & Centre observation & Candidats\\
\midrule
S & Mélange $C_\gamma$, $7\times20$ couples & $A\mu$ & 140\\
M0 & Covariance signal $C_{\rm emp}$, $\gamma=0$ & $A\mu$ & 20\\
Diag & $\operatorname{diag}(\operatorname{diag}C_{\rm emp})$, cosinus & $A\mu$ & 20\\
E & $C_{cy}(C_{yy}+\lambda t_yI)^{-1}$ empirique & $\mu_y$ & 20\\
H1 & Tikhonov historique, pénalité $Q/Q_{\max}$ & $\mu_y$ & 20\\
H1m & Extrémité $\gamma=1$, grille $\alpha$ appariée & $A\mu$ & 20\\
\bottomrule
\end{tabular}
\end{center}
\par\vspace{2mm}
\begin{block}{Deux distinctions nécessaires}
$\gamma=0$ n'est pas en général le LMMSE empirique E : blocs de covariance et centrage diffèrent. H1m et H1 peuvent aussi différer par leur centre et leur grille.
\end{block}
\scriptsize $\gamma\in\{0;0,01;0,03;0,1;0,3;0,6;1\}$ ; $\alpha\in\operatorname{logspace}(-3,3,20)$. Le supplément de recherche de S est assumé ; aucune égalité de budget n'est revendiquée.
\end{frame}

\begin{frame}{Règle figée avant le test final}
\begin{columns}[T]
\column{.49\textwidth}
\textbf{Données et sélection}
\begin{itemize}
\item $N=32,64,256$, cinq entraînements indépendants ; tailles imbriquées.
\item 256 objets de validation nominale, deux bruits ; réglage supervisé uniquement sur validation.
\item 400 objets test, quatre bruits appariés ; perte quadratique L$^2$ sur les 128 modes.
\end{itemize}
\column{.49\textwidth}
\textbf{Un contraste : S/M0, $N=64$}
\begin{itemize}
\item Gain utile avec 64 capteurs : ratio de risque $\leq0,97$.
\item Protection avec 32 contigus : ratio de risque $\leq1,03$.
\item 4\,000 bootstraps par objet ; bornes 99,5\,\%, dix inégalités simultanées par Bonferroni.
\end{itemize}
\end{columns}
\par\vspace{3mm}
\small Gel : \texttt{8c23fcc}. Seuils conventionnels fixés avant les nouveaux résultats. Incertitude conditionnelle ; couverture bootstrap approximative, sans intervalle populationnel sur l'aléa d'entraînement.
\end{frame}

\begin{frame}{Contraste principal : la règle complète n'est pas satisfaite}
\centering\plot{primary}{.98}\figsource{primary}
\par\vspace{2mm}
\small Les bornes du scénario principal dépassent toutes 0,97. Les cinq bornes de protection restent sous 1,03. Les deux exigences doivent être réunies dans chaque entraînement.
\end{frame}

\begin{frame}{Cinq entraînements : les valeurs derrière la décision}
\begin{center}
\begin{tabular}{rrrrr}
\toprule
Jeu & Ratio principal & Borne sup. & Ratio partiel & Borne sup.\\
\midrule
0 & 0,9607 & 0,9844 & 1,0000 & 1,0000\\
1 & 0,9759 & 1,0001 & 1,0060 & 1,0291\\
2 & 0,9707 & 0,9903 & 1,0004 & 1,0158\\
3 & 0,9823 & 1,0097 & 1,0037 & 1,0074\\
4 & 0,9709 & 0,9858 & 1,0126 & 1,0267\\
\bottomrule
\end{tabular}
\end{center}
\par\vspace{3mm}
\begin{itemize}
\item Principal : amélioration ponctuelle dans les cinq jeux, mais gain utile de 3\,\% non établi.
\item Partiel : quatre ratios ponctuels dépassent 1 ; ils demeurent compatibles avec la marge préfixée de 3\,\%.
\item Une absence de gain utile établi n'est pas une preuve d'équivalence.
\end{itemize}
\par\vspace{2mm}
{\tiny\src{results/analysis/report.json}{Ratios et bornes exacts}\,|\,\src{results/analysis/bootstrap.npz}{Tirages bootstrap}\,|\,\src{src/analysis.py}{Analyse}}
\end{frame}

\begin{frame}{Tailles d'apprentissage : des performances dépendantes du régime}
\centering\plot{training}{1}\figsource{training}
\par\vspace{2mm}
\small Les tailles sont imbriquées dans chaque entraînement ; le test reste commun. La dispersion entre cinq entraînements est distincte de l'incertitude conditionnelle. Comparaisons secondaires descriptives.
\end{frame}

\begin{frame}{Le poids de la cible H1 dépend de la couverture}
\centering\plot{gamma}{.98}\figsource{gamma}
\par\vspace{2mm}
\small À $N=64$, capteurs complets : $\gamma=(1;1;0,6;1;0,6)$. Avec 32 contigus : $\gamma=(0;0,1;0,1;0,03;0,1)$. Une sélection à l'extrémité est conservée, sans élargir la grille après test.
\end{frame}

\begin{frame}{Même $\alpha$ : isoler l'effet du changement de covariance}
\centering\plot{matched}{.98}\figsource{matched}
\par\vspace{2mm}
\small M0\_at\_S et H1m\_at\_S utilisent le $\alpha$ choisi pour S : seuls les mélanges de covariance changent. Ce diagnostic n'est pas une comparaison de trois méthodes chacune réglée au mieux, ni un second contraste confirmatoire.
\end{frame}

\begin{frame}{Robustesse : les dégradations restent dans le bilan}
\begin{columns}[c]
\column{.72\textwidth}
\centering\plot{robustness}{1}\figsource{robustness}
\column{.26\textwidth}
\small S face à M0, $N=64$ :
\par\vspace{3mm}
Rugueux, complet :\par
\alert{+1,37\,\%} de risque.
\par\vspace{3mm}
Hauteur $+10\,\%$, complet :\par
\alert{+2,66\,\%}.
\par\vspace{3mm}
Hauteur $+10\,\%$, partiel :\par
\alert{+1,89\,\%}.
\par\vspace{3mm}
\scriptsize Ratios des risques moyens sur cinq entraînements ; descriptifs, sans preuve secondaire corrigée.
\end{columns}
\par\vspace{2mm}
\small Au nominal complet, Diag obtient aussi un risque moyen inférieur à S. Aucun avantage universel de la covariance mélangée n'est établi.
\end{frame}

\begin{frame}{Contrôles ciblés et limites de portée}
\begin{columns}[T]
\column{.49\textwidth}
\textbf{Vérifications}
\begin{itemize}
\item Audit des archives de l'étude 2 : gel, partitions, scalaires et projections ; aucune anomalie substantielle détectée.
\item 17 tests réussis : extrémités, covariance singulière, E/H1 historiques, sélection et bootstrap.
\item Pilote séparé réussi ; aucune panne numérique dans la campagne.
\end{itemize}
\column{.49\textwidth}
\textbf{Limites}
\begin{itemize}
\item Simulation gaussienne, même géométrie ; aucune donnée expérimentale réelle.
\item Cible H1 imposée, grilles finies et validation uniquement nominale.
\item Cinq entraînements ; incertitude conditionnelle aux modèles observés.
\item S examine 140 candidats, contre 20 par comparateur.
\end{itemize}
\end{columns}
\par\vspace{3mm}
\scriptsize Campagne seule : 7,69 s CPU, pic 98,1 Mo RAM, sorties 63,5 Mo. Plafonds globaux : 4 h CPU, 2 Gio RAM, 200 Mo de nouvelles sorties. Aucun service payant, GPU ou FreeFem.\par
\src{AUDIT_STUDY2.md}{Audit préalable}\,|\,\src{results/tests.json}{Tests ciblés}\,|\,\src{results/final/report.json}{Bilan de campagne}
\end{frame}

\begin{frame}{Conclusion, attribution et contact}
\begin{block}{Conclusion}
La cible H1 apporte une amélioration ponctuelle au nominal complet ; le gain utile fixé à 3\,\% n'est pas établi. La marge de non-dégradation partielle est satisfaite conditionnellement, mais des dégradations secondaires subsistent.
\end{block}
\small
Antériorités : Ledoit--Wolf (2004), Alberti et al. (NeurIPS 2021), Burger--Kabri (2024), Tajani--Abouchabaka (2018). \src{SOURCES.md}{Références et portée des lectures}. Aucun algorithme nouveau revendiqué ; aucune garantie transférée sans preuve.
\par\vspace{2mm}
Code, analyse et rédaction avec assistance d'IA ; validation académique relevant de l'auteur et de l'encadrant. Les logos ne valent pas approbation des résultats.
\par\vspace{2mm}
\textbf{Halhoul Mhamed}\par
\href{mailto:halhoul.mhamed@etu.uae.ac.ma}{halhoul.mhamed@etu.uae.ac.ma}\par
\href{mailto:halhoulmhamed@gmail.com}{halhoulmhamed@gmail.com}\par
\href{https://github.com/halhoulmhamed-droid/inverse-laplace-learning}{Dépôt public inverse-laplace-learning}\par
\par\vspace{2mm}
{\tiny Version scientifique : \texttt{\DataCommit}}
\end{frame}
\end{document}
